% Arquitectura del modelo integrado del motor VCR (NB10b): que resuelve cada submodelo y como se acoplan.
\documentclass[tikz,border=4mm]{standalone}
\usepackage{amsmath}
\usetikzlibrary{arrows.meta,positioning,fit,backgrounds}
\begin{document}
% Estilo comun de los diagramas TikZ del modulo VCR (NB10b) -- paleta del proyecto
\definecolor{manivela}{HTML}{2A78D6}
\definecolor{inferior}{HTML}{4A3AA7}
\definecolor{control}{HTML}{1BAF7A}
\definecolor{superior}{HTML}{EB6834}
\definecolor{tinta}{HTML}{3A3A38}
\definecolor{metal}{HTML}{C3C2B7}
\definecolor{perdida}{HTML}{E34948}
\definecolor{fuerza}{HTML}{C0392B}
\tikzset{
  barra/.style={line width=3.2pt, line cap=round},
  pin/.style={circle, fill=white, draw=tinta, line width=0.9pt, inner sep=1.6pt},
  fijo/.style={circle, fill=tinta, inner sep=2.6pt},
  vec/.style={-{Stealth[length=2.6mm]}, line width=0.9pt},
  F/.style={-{Stealth[length=3mm]}, line width=1.2pt, fuerza},
  cota/.style={{Stealth[length=1.8mm]}-{Stealth[length=1.8mm]}, line width=0.5pt, tinta},
  lbl/.style={font=\small, text=tinta},
}
\newcommand{\tierra}[1]{\fill[pattern=north east lines, pattern color=tinta] ($(#1)+(-5,-3)$) rectangle ($(#1)+(5,-6)$);
  \draw[tinta, line width=0.8pt] ($(#1)+(-5,-3)$) -- ($(#1)+(5,-3)$);}

\begin{tikzpicture}[
  caja/.style={draw=#1, line width=1.1pt, rounded corners=3pt, fill=#1!8, align=center, font=\small, text width=4.3cm, minimum height=1.5cm, inner sep=4pt},
  flecha/.style={-{Stealth[length=2.6mm]}, line width=1pt, tinta},
  node distance=0.9cm and 1.3cm]
  \node[caja=manivela] (cin) {\textbf{Cinem\'atica 2D}\\ lazos + jacobiano\\ $y_D(\theta),\ \dot y_D,\ \ddot y_D$};
  \node[caja=control, left=of cin] (vcr) {\textbf{Control VCR}\\ $R_{OE}\ \leftrightarrow\ r_c\in[8,14]$\\ (eje excéntrico $E$)};
  \node[caja=inferior, below=of cin] (vol) {\textbf{Volumen}\\ $V=V_c+A_p(y_{\text{culata}}-y_D)$\\ $\dot V=-A_p\,\dot y_D$};
  \node[caja=superior, below=of vol] (termo) {\textbf{Termodin\'amica 0D, 2 zonas}\\ 1.\textsuperscript{a} ley en $t$; Woschni\\ $p(t),\,T_u(t),\,T_b(t)$};
  \node[caja=perdida, right=of termo] (quim) {\textbf{Cin\'etica qu\'imica detallada}\\ Cantera, iso-octano\\ 143 especies / 1178 reacc.\\ autoencendido del gas final};
  \node[caja=perdida, left=of termo] (wiebe) {\textbf{Llama (chispa)}\\ $\dot x_b(t)$ tipo Wiebe\\ CA50 limitado por\\ detonaci\'on (KLSA)};
  \node[caja=manivela, below=of termo] (ne) {\textbf{Din\'amica (Newton--Euler)}\\ $\mathbf A(\theta)\mathbf x=\mathbf b$, 14 inc\'ognitas\\ fuerzas en pasadores, $T_O(\theta)$\\ {\footnotesize (con $\vec a_{cm},\alpha_i$ de la cinem\'atica)}};
  \node[caja=tinta, right=of ne] (des) {\textbf{Desempe\~no}\\ IMEP, FMEP, BMEP\\ par y potencia vs. rpm};
  \draw[flecha] (vcr) -- (cin);
  \draw[flecha] (cin) -- (vol);
  \draw[flecha] (vol) -- node[right, font=\footnotesize] {$V(t),\dot V(t)$} (termo);
  \draw[flecha] (quim.west) -- node[above, font=\footnotesize] {$\dot\omega_k$} (termo.east);
  \draw[flecha] (termo.north east) to[bend left=25] node[right, font=\footnotesize] {$p,T_u$} (quim.north);
  \draw[flecha] (wiebe) -- node[above, font=\footnotesize] {$\dot m_{u\to b}$} (termo);
  \draw[flecha] (termo) -- node[right, font=\footnotesize] {$F_{gas}=(p-p_c)A_p$} (ne);
  \draw[flecha] (ne) -- (des);
  \draw[flecha] (termo.south east) to[bend left=10] (des.north);
  \node[font=\footnotesize, text=tinta, align=left, anchor=north west] at ([yshift=-0.4cm]wiebe.south west) {cuasi-est\'atico: $\omega$ = cte.\\ (verificado con el volante, \S8)};
\end{tikzpicture}
\end{document}
